\documentclass[sigplan,twocolumn,10pt]{acmart}

\renewcommand\footnotetextcopyrightpermission[1]{}
\settopmatter{printfolios=true,printacmref=false}
\pagestyle{plain}

\def\red#1{\textcolor{red}{#1}}


\AtBeginDocument{%
  \providecommand\BibTeX{{%
    Bib\TeX}}}

%% Rights management information.  This information is sent to you
%% when you complete the rights form.  These commands have SAMPLE
%% values in them; it is your responsibility as an author to replace
%% the commands and values with those provided to you when you
%% complete the rights form.
\setcopyright{acmlicensed}
\copyrightyear{2018}
\acmYear{2018}
\acmDOI{XXXXXXX.XXXXXXX}
%% These commands are for a PROCEEDINGS abstract or paper.
\acmConference[Conference acronym 'XX]{Make sure to enter the correct
  conference title from your rights confirmation email}{June 03--05,
  2018}{Woodstock, NY}
%%
%%  Uncomment \acmBooktitle if the title of the proceedings is different
%%  from ``Proceedings of ...''!
%%
%%\acmBooktitle{Woodstock '18: ACM Symposium on Neural Gaze Detection,
%%  June 03--05, 2018, Woodstock, NY}
\acmISBN{978-1-4503-XXXX-X/2018/06}


%%
%% Submission ID.
%% Use this when submitting an article to a sponsored event. You'll
%% receive a unique submission ID from the organizers
%% of the event, and this ID should be used as the parameter to this command.
%%\acmSubmissionID{123-A56-BU3}

%%
%% For managing citations, it is recommended to use bibliography
%% files in BibTeX format.
%%
%% You can then either use BibTeX with the ACM-Reference-Format style,
%% or BibLaTeX with the acmnumeric or acmauthoryear sytles, that include
%% support for advanced citation of software artefact from the
%% biblatex-software package, also separately available on CTAN.
%%
%% Look at the sample-*-biblatex.tex files for templates showcasing
%% the biblatex styles.
%%

%%
%% The majority of ACM publications use numbered citations and
%% references.  The command \citestyle{authoryear} switches to the
%% "author year" style.
%%
%% If you are preparing content for an event
%% sponsored by ACM SIGGRAPH, you must use the "author year" style of
%% citations and references.
%% Uncommenting
%% the next command will enable that style.
%%\citestyle{acmauthoryear}

\settopmatter{authorsperrow=4}


\begin{document}

\title[Advanced Computing Project (2025/26)]{Storage system for optimizing LLM checkpointing in HPC}

\author{André Miranda}
\email{pg60238@alunos.uminho.pt}
\affiliation{
  \institution{Universidade do Minho}
  \city{Braga}
  \country{Portugal}
}

\author{Gustavo Barros}
\email{pg61527@alunos.uminho.pt}
\affiliation{
  \institution{Universidade do Minho}
  \city{Braga}
  \country{Portugal}
}

\author{José Soares}
\email{pg60277@alunos.uminho.pt}
\affiliation{
  \institution{Universidade do Minho}
  \city{Braga}
  \country{Portugal}
}

\author{Miguel Carvalho}
\email{pg61153@alunos.uminho.pt}
\affiliation{
  \institution{Universidade do Minho}
  \city{Braga}
  \country{Portugal}
}

\renewcommand{\shortauthors}{Miranda et al.}

\begin{abstract}
  Training Large Language Models (LLMs) in High-Performance Computing (HPC) environments is a time-consuming process that requires frequent 
  checkpointing to ensure fault tolerance. However, concurrent writes of massive model states often saturate the shared Parallel File System 
  (PFS), leading to significant I/O contention and training stalls. This paper proposes a policy-driven, two-tier storage orchestration system
  designed to mitigate the "noisy neighbor" problem in HPC clusters. By transparently redirecting checkpoints to node-local SSDs and employing
  a centralized orchestrator to manage asynchronous migration to the PFS via a Token Bucket algorithm, we enable granular Quality of Service (QoS)
  control. Our evaluation on the Deucalion supercomputer demonstrates that ... 
\end{abstract}

\keywords{High-Performance Computing, Deucalion, A64FX, AMD EPYC 7742, Memory Hierarchy, Large Language Models, Checkpointing, Parallel File Systems, 
I/O Contention, Quality of Service (QoS), Distributed Systems, Resource Management..}


\received{25 May 2026}

\maketitle


\section{Introduction}

The rapid growth of Large Language Models (LLMs) has significantly increased the computational and storage demands of modern 
High-Performance Computing (HPC) systems. Training such models requires long-running, distributed jobs over clusters of GPUs 
or accelerators, where fault tolerance is essential to avoid losing weeks of computation due to hardware failures, network 
instability, or software errors. Checkpointing—the periodic saving of model states, optimizer parameters, and training 
metadata—has emerged as the de facto mechanism to ensure progress is preserved. However, as model sizes scale to billions or 
even trillions of parameters, checkpoint operations now generate massive I/O workloads, often exceeding tens to hundreds of 
gigabytes per save.

In shared HPC environments, these I/O-intensive workloads are typically directed to a Parallel File System (PFS), such as 
Lustre or GPFS, which is concurrently accessed by hundreds of users and applications. When multiple deep learning jobs attempt 
to write checkpoints simultaneously—as is common at the end of training epochs—the resulting synchronized I/O bursts can saturate 
the PFS, leading to severe contention. This phenomenon, known as the Noisy Neighbor Problem, degrades performance cluster-wide, 
increasing latency, reducing effective throughput, and even stalling training processes entirely. Such stalls not only waste 
computational resources but also disrupt the workflows of unrelated jobs sharing the same filesystem, creating a cascading effect 
that undermines the efficiency of the entire HPC ecosystem.

The central insight driving this work is that checkpointing is not latency-critical for persistence, but it is latency-critical 
for training continuation. In other words, training can resume as soon as the checkpoint is safely written to fast, temporary 
storage (e.g., node-local SSDs), while the durable persistence to the PFS can occur asynchronously in the background. This decoupling 
of checkpoint generation from long-term storage unlocks the potential to smooth I/O bursts, reduce contention, and improve overall system 
stability without sacrificing fault tolerance.

To address this challenge, we propose a policy-driven, two-tier storage orchestration system that combines node-local SSD staging 
with centralized rate control. The system is designed to transparently intercept checkpoint writes, redirect them to fast local 
storage, and asynchronously migrate them to the PFS under strict Quality of Service (QoS) constraints. By doing so, we minimize 
the stall time experienced by training jobs while regulating global I/O pressure to prevent PFS saturation.

\section{Background and Motivation}

\subsection{Checkpointing in HPC}

Checkpointing is a fault-tolerance mechanism in distributed deep learning that periodically saves the model state, including 
parameters, optimizer states, and gradients, to stable storage. As models scale to billions of parameters, checkpoint sizes 
can reach tens to hundreds of gigabytes, making I/O operations a critical bottleneck. Since large-scale training jobs may execute 
for hours or even weeks, failures caused by hardware faults, network instability, or software errors can lead to significant 
computational losses. To mitigate this problem, applications periodically save their execution state to stable storage, enabling 
recovery from the most recent checkpoint instead of restarting from the beginning.


\subsection{Parallel File Systems and Contention}


Parallel file systems are the storage backbone of HPC infrastructures, designed to provide massive data throughput for thousands of 
concurrent jobs. In this work, we focus on Lustre-like file systems, which are widely deployed across TOP500 supercomputers. A typical 
Lustre architecture decouples metadata from file data: Metadata Servers (MDSs) manage the namespace, while Object Storage Servers (OSSs) 
and their underlying Object Storage Targets (OSTs) handle the heavy lifting of data read and write operations. To achieve high bandwidth, 
large files are typically striped across multiple OSTs, allowing compute nodes to write data in parallel via kernel-level clients using 
standard POSIX interfaces.

While this architecture provides massive aggregate bandwidth, it is increasingly challenged by the I/O patterns of modern Deep Learning 
(DL) workloads, particularly due to model checkpointing. As neural networks grow exponentially in size, saving a model's state routinely 
generates monolithic files or large shards ranging from tens to hundreds of gigabytes. Unlike dataset loading, which is continuous and 
read-intensive, checkpointing introduces periodic, extremely write-intensive I/O bursts.

When multiple large-scale distributed training jobs attempt to write these massive checkpoints to the PFS simultaneously they can saturate 
the OSSs and the cluster's network interconnect. This severe I/O contention not only degrades the overall storage performance for all users 
in the HPC cluster but also severely impacts the training efficiency itself. Because standard checkpointing mechanisms are typically synchronous, 
expensive GPU compute nodes are forced to stall, wasting valuable computational resources while waiting for the bottlenecked storage system to 
complete the massive write operations.

\subsection{Quality of Service in HPC Storage}

To mitigate this severe I/O contention, Quality of Service (QoS) mechanisms are increasingly necessary in HPC environments. QoS refers to the 
ability to control and prioritize resource allocation—such as storage bandwidth—to ensure fairness, predictability, and overall cluster stability. 
In the context of parallel file systems, QoS can be enforced dynamically to prevent concurrent deep learning jobs from monopolizing shared resources 
during their massive, synchronized checkpointing phases.

\subsection{Token Bucket Algorithm}

A foundational technique for enforcing such bandwidth QoS is the Token Bucket algorithm, \red{TODO}


\subsection{Motivation}

The main motivation of this work stems from the severe I/O bottlenecks observed when training large-scale Deep Learning (DL) models on HPC clusters. 
As model parameters scale into the billions, their corresponding checkpoints grow to tens or hundreds of gigabytes. Currently, when distributed training 
jobs attempt to write these massive checkpoints to a shared Parallel File System (PFS) simultaneously.

This uncoordinated burst of I/O traffic instantly saturates the storage network, causing extreme contention. Consequently, expensive compute nodes are 
forced to stall while waiting for synchronous write operations to finish, wasting valuable computational time and degrading the storage performance for the 
entire cluster. This work is motivated by the critical need to eliminate these stalls. By developing a dynamic, QoS-aware orchestration layer, we aim to decouple 
checkpoint generation from PFS flushing, ensuring smooth bandwidth utilization and maximizing training efficiency.



\section{System Design}

The proposed system adopts a two-tier storage architecture combined with a centralized orchestration mechanism to regulate checkpoint 
traffic in shared HPC environments.

The main design goals are:
\begin{itemize}
    \item Reduce checkpoint latency perceived by training applications
    \item Minimize congestion on the shared Parallel File System (PFS)
    \item Provide fairness and QoS isolation between concurrent jobs
    \item Preserve transparency and compatibility with existing HPC workflows
\end{itemize}

Instead of writing checkpoints directly to the PFS, the system temporarily stages data in node-local SSDs and asynchronously migrates 
it under centralized control. This decouples checkpoint generation from persistent storage operations, reducing synchronized I/O bursts.


\subsection{Design Overview}

The system consists of three main components:
\begin{itemize}
    \item \textbf{Interception Layer}: Transparently redirects checkpoint writes to local storage.
    \item \textbf{Local Storage Tier}: Node-local SSDs for fast staging.
    \item \textbf{Control Plane}: Central orchestrator managing global I/O behavior and enforcing QoS policies.
\end{itemize}

When a checkpoint operation is triggered, the application writes data to the local SSD instead of directly accessing the PFS. Once the 
local write completes, the application can immediately resume execution. Checkpoints are later transferred to persistent storage through 
controlled background migration managed by the control plane.

This design minimizes training stalls while smoothing I/O pressure across the cluster.

\subsection{Local Storage Tier}

Each compute node provides local SSD storage used as a temporary checkpoint staging area. Compared to remote PFS access, local SSDs offer 
significantly lower latency and higher write throughput.

This storage tier absorbs bursty checkpoint traffic and enables rapid checkpoint completion, allowing training processes to continue execution 
without waiting for persistence to complete.

In addition to storing checkpoint data, the local tier maintains metadata describing:
\begin{itemize}
    \item Checkpoint identifiers
    \item Migration status
    \item File integrity information
    \item Persistence progress
\end{itemize}

The local storage layer therefore acts as a buffering mechanism between high-frequency checkpoint generation and comparatively slower shared 
storage operations.

\subsection{Control Plane}

The control plane is responsible for globally coordinating checkpoint migration across compute nodes. It maintains a system-wide view of:
\begin{itemize}
    \item Active jobs
    \item Pending checkpoint queues
    \item PFS utilization
    \item Allocated bandwidth quotas
\end{itemize}

Rather than allowing unrestricted checkpoint flushing, the control plane regulates migration traffic using a Token Bucket-based QoS mechanism. 
Each job receives a configurable bandwidth allocation that determines the rate at which checkpoint data may be transferred to the PFS.

This centralized coordination prevents individual jobs from monopolizing storage bandwidth and mitigates the noisy neighbor effect commonly 
observed in shared HPC systems.

By smoothing bursty migration traffic over time, the control plane improves overall storage stability and fairness while maintaining predictable 
checkpoint behavior.



\section{Implementation}

\subsection{Token Bucket Model}

To regulate PFS access, each job is assigned a token bucket defining its allowed write rate. Tokens accumulate over time and are consumed during 
data transfers. \red{TODO: Describe the specific parameters of the token bucket and how they are dynamically adjusted based on system conditions.}

\subsection{Scheduling Policy}

The orchestrator schedules \red{TODO: Describe the scheduling policies for checkpoint migration.}

\subsection{Prototype Design}

\red{TODO: Describe the prototype implementation, including the interception mechanism, local storage management, and control plane architecture.}

\subsection{Data Management}

Checkpoint files are stored locally with metadata tracking their migration status. A background daemon handles transfers to the PFS.



\section{Evaluation}

Experiments were conducted on the Deucalion supercomputer, featuring heterogeneous compute nodes and a shared PFS. Workloads simulate LLM checkpointing 
behavior with varying sizes and frequencies.

\subsection{Workloads}
We evaluate the system using synthetic and realistic checkpointing workloads designed to mimic Large Language Model (LLM) training behavior:
\begin{itemize}
    \item TODO
\end{itemize}

\subsection{Metrics}

\begin{itemize}
    \item Checkpoint latency
    \item Training stall duration
    \item PFS throughput and variability
    \item Fairness index across jobs
    \item I/O congestion
    \item CPU/GPU and memory usage
\end{itemize}




\section{Related Work}

The optimization of checkpointing in High-Performance Computing (HPC) has been extensively studied, primarily focusing on reducing write 
latency and mitigating Parallel File System (PFS) congestion.

\subsection{Architectural and Tiering Approaches}
Several systems utilize hierarchical storage to hide I/O latency. \textbf{Check-N-Run} \cite{CheckNRun} focuses on staging checkpoints in local node resources 
(DRAM or SSDs), allowing training to resume immediately while data is moved to the PFS in the background. However, it lacks a global coordination 
mechanism to prevent PFS saturation during the background flush. \textbf{DeepFreeze} \cite{DeepFreeze} (similar to Monarch) employs transparent storage tiering to 
accelerate the write path by abstracting complex local tiers like NVMe or PMEM from the user. While effective at the node level, these 
systems often ignore the "noisy neighbor" effect in shared environments.

\subsection{Adaptive and Optimization Techniques}
Other research focuses on the frequency and size of checkpoints. \textbf{CheckFreq} \cite{CheckFreq} introduces an algorithmic approach to identify optimal checkpoint
boundaries and pipelines the process with training iterations, dynamically adjusting frequency to maintain overhead below a threshold. 
\textbf{BitSnap} \cite{BitSnap} addresses the data volume by employing sparsification and quantization, reducing checkpoint sizes from 16-bit to 4-bit weights before
they reach the I/O layer. While these reduce total I/O pressure, they do not manage the storage hierarchy or provide holistic Quality of Service
(QoS) guarantees.

\subsection{Software-Defined Storage and QoS Control}
Our work is heavily informed by \textbf{PADLL} \cite{PADLL}, a storage middleware that enables application-agnostic QoS control for both data and metadata workflows.
PADLL utilizes a decoupled Software-Defined Storage (SDS) architecture, where data plane stages intercept POSIX calls and a centralized 
control plane coordinates I/O rates across all jobs. While PADLL is a general-purpose QoS framework, our system specializes this paradigm for
the specific lifecycle of LLM checkpoints—integrating node-local SSD staging with PADLL-inspired Token Bucket rate limiting to ensure cluster harmony 
during massive model state migrations.




\section{Conclusions}




\section{Future Work}

Future work will focus on...

\bibliographystyle{ACM-Reference-Format}
\bibliography{references}

\end{document}
\endinput
%%
%% End of file `sample-sigplan.tex'.
